\begin{problem}[Specialization geometry of the crossing axes]
Let
\[
R=k[x,y]/(xy).
\]
Show that every prime of $R$ contains $\bar x$ or $\bar y$.  Identify the two
component-generic points and, when $k$ is algebraically closed, describe the closed points
on the two components and their specialization relations.
\end{problem}

\paragraph{Why this problem matters.}
This is the basic reducible example in which one closed point can be a specialization of
more than one generic point.

\paragraph{Useful ideas before starting.}
A prime containing $\bar x\bar y=0$ must contain one factor.  Then reduce to the spectra of
$k[x]$ and $k[y]$.

\paragraph{Guided hints.}
Primes containing $\bar x$ correspond to primes of
\[
R/(\bar x)\cong k[y].
\]

\begin{solution}
For any prime $\mathfrak p\subseteq R$,
\[
\bar x\bar y=0\in\mathfrak p.
\]
Primality gives
\[
\bar x\in\mathfrak p
\quad\text{or}\quad
\bar y\in\mathfrak p.
\]
Thus every point lies on at least one of the two components.

The ideals
\[
(\bar x),\qquad(\bar y)
\]
are prime because their quotients are $k[y]$ and $k[x]$.  They are minimal primes, hence
the generic points of the two irreducible components.

If $k$ is algebraically closed, the closed points on $V(\bar x)$ are
\[
(\bar x,\bar y-a),\qquad a\in k,
\]
and on $V(\bar y)$ they are
\[
(\bar x-a,\bar y),\qquad a\in k.
\]
At $a=0$ these two descriptions coincide in the crossing point
\[
(\bar x,\bar y).
\]
Each closed point on an axis is a specialization of that axis's generic point, and the
crossing point is a specialization of both.
\end{solution}

\paragraph{What happened mathematically.}
The specialization order records incidence: the common closed point lies in the closures
of both component-generic points.

\paragraph{Extensions.}
Repeat the analysis for $k[x,y]/(xy(x-1))$.

\paragraph{Provenance.}
New pointwise continuation of the legacy crossing-axes examples already used in VI/05 and
VI/08.
